\section{Discussion}
\label{sec:discussion}

\subsection{Why interpreted updates only}
The survey's strongest empirical result is not that self-evolution fails; it is that the field has not yet established when it works. Across 93 primary comparisons, none meet all four evaluation claims, and Activation and Attribution are rarely fully evidenced \citep{wei2026harnessrsi}. For an academic harness, the asymmetry is worse than in coding benchmarks. A code patch that weakens a check can make every later manuscript look compliant, and the damage may not appear until a retraction, a failed replication, or a reviewer objection. The survey's failure-mode literature makes the same point from several directions: harness tampering, phantom guardrails, self-authored verification, and memory reward inflation are not hypothetical \citep{wang2026tampering,wang2026phantom,guo2026selfverification,asadolahi2026inflation}.

\system{} responds by making the smallest update that can persist: a short lesson. The deterministic audit, the vendored engine, the project ledger, the plugin code, and the admission policy remain outside the loop. This is not a claim that lesson learning is safe in an absolute sense. It is a claim about blast radius. A bad lesson can waste a turn or mislead a draft; it cannot rewrite the check that would catch the error, alter a fact, or approve a waiver. The author can see it, forget it, or purge the store. The normal release process remains the only path to executable change.

\subsection{Autonomy without approval}
The system admits lessons without author approval. This is a deliberate departure from the author-only pattern used for waivers, fact proposals, plans, and disclosures in the writing engine. The distinction is authority. Approving a fact or waiving a finding changes the project's evidentiary state; a lesson changes only a piece of advice injected into a later turn. The former must remain an author decision; the latter can be automated if it is bounded and reversible.

Autonomy is still constrained. A diagnosis lesson starts as a trial with a randomized control; an active lesson keeps a control sample; caps limit the library; a time-to-live removes stale advice; the lint rejects a set of dangerous instructions; a contradiction check prevents a new lesson from silently reversing an old one; and explicit preferences require a durable cue and cannot be created by the model. The result is not ``the agent can do anything.'' It is ``the agent can propose and test advice, and the user can inspect and undo it.''

The autonomy/control trade-off is the main design question a reader should challenge. A stricter system would require approval for every lesson, which would likely reduce usage to near zero. A looser system would let the agent edit skills, prompt sections, or checks, which the evidence does not support. The middle ground is a lesson store with a visible lifecycle and a fixed oracle.

\subsection{Relationship to prior self-improving systems}
\system{} is closest to the experience-mining and gated-consolidation lines. Reflexion and ExpeL store reflections or cross-task insights \citep{shinn2023reflexion,zhao2024expel}; ACE grows a context playbook through deltas \citep{zhang2026agentic}; ReasoningBank injects consolidated reasoning strategies \citep{ouyang2026reasoningbank}; SkillOpt, SkillClaw, and VALVE gate skill revisions on validation \citep{yang2026skillopt,ma2026skillclaw,wang2026valve}. Our lessons are smaller than a skill and are admitted online rather than through a centralized validation panel. The trial/control lifecycle is a local, low-cost analogue of those gates.

The system is farther from prompt and workflow optimizers such as APE, OPRO, TextGrad, GEPA, DSPy, and AFlow \citep{zhou2023large,yang2024large,yuksekgonul2025textgrad,agrawal2026gepa,khattab2024dspy,zhang2025aflow}. Those methods optimize a program against a dataset and a metric. \system{} has a non-stationary task stream, a heuristic outcome, and no offline optimization budget. It therefore uses a remedy catalogue and online evidence rather than a search over prompts or graphs.

It is farthest from repository-level self-improvement such as SICA, DGM, HGM, DarwinX, MOSS, Ouroboros, and AgentDevel \citep{robeyns2025selfimproving,zhang2026darwin,wang2025huxley,zhang2026darwinx,cai2026moss,razzhigaev2026ouroboros,zhang2026agentdevel}, and from component-patching systems such as Self-Harness, HarnessCompass, HarnessFix, HarnessX, and HarnessBank \citep{zhang2026selfharness,zhang2026harnesscompass,chen2026harnessfix,chen2026harnessx,luo2026harnessbank}. Those systems explore a much larger space and can produce larger gains, but they also require a code-review or deployment boundary. \system{} keeps that boundary and learns only the interpreted layer.

\subsection{The oracle is strong but narrow}
The deterministic writing oracle is the main reason the loop can run without an author in the admission path. It supplies stable finding identifiers, a gate, report hashes, and a ledger. Those are excellent signals for lesson diagnosis: a numeric finding, a missing citation key, or a compile error is a real event with a stable name. The oracle is also narrow. It cannot tell whether an argument is original, whether a statistical model is appropriate, or whether a citation supports the claim in the way the prose implies. The system therefore treats the oracle as a detector, not as a measure of research quality.

The narrowness has a second consequence: the loop can be gamed by doing less. An agent could avoid new findings by editing little, or by removing the number that triggered a check. The outcome score partially addresses this by including correction markers, gate continues, blocking findings, and a success threshold, but it does not include coverage or task completion. This is an open problem. A future version should add brief and inventory coverage, reviewer-item resolution, and unnecessary-handoff rate to the outcome vector. The current implementation records enough state to add them without changing the store format.

\subsection{What the corpus evidence adds}
\label{sec:corpus-discussion}
The nine-document experiment is the strongest real-data evidence in the paper, and it changes the discussion in three ways. First, the signal stream is real and concentrated: 1,158 findings on 690 units, dominated by numeric and bibliography checks. A deployed project with a populated source registry and bound facts would not have the same mix; the corpus is a cold-start view. This matters because the lesson catalogue is calibrated for the same check families, so coverage on a mature project could be lower or higher. The experiment should be repeated on projects that have already been curated, where the findings reflect manuscript problems rather than missing setup.

Second, defect injection is not the same as lesson efficacy. All 31 injected defects were detected, but detection is the oracle doing its job. The transfer result is more relevant to the RSI claim: leave-one-project-out lessons covered 23/23 specific families, and pooling proposals raised a sparse family from 2/3 to 3/3. That is evidence that a shared local lesson store can address the cold-start problem for an individual researcher, at a cost of 62--82 tokens per document. It is not evidence that the agent follows the lesson or that the manuscript improves. The next study must measure adherence and outcomes in a live session.

Third, the negative results identify the real bottlenecks. Task routing is 56.7\% accurate on a hand-labeled set, and pattern-based safety linting drops to 0.53 F1 on paraphrased unsafe lessons. The learning policy is not the weak link; the interfaces around it are. A local semantic classifier and a semantic safety filter are the obvious next investments, but they must not become editable by the loop. The fixed oracle and the text-only mutable surface are what keep the system's failure modes bounded even when the lint misses a paraphrase.

\subsection{Local-first and the remote option}
Research data and manuscripts are often confidential. A local store is therefore the default, and the loop makes no network calls unless the user opts into the remote judge. The remote judge is a veto, not an optimizer; it receives lesson text, task class, action, and evidence counts, never the manuscript or the error signature. This is a compromise between privacy and the value of an independent check. It is not a substitute for a proper external evaluation, which would require consent, de-identification, and a data-processing agreement.

A local-first design also has costs. Lessons are scoped to one user's machine and one user's workflow; they do not benefit from a population of researchers. Cross-project transfer is assumed by the user-level store but not measured. A lesson learned in one manuscript may be irrelevant or wrong in another. The task-class scope reduces the risk, and the injection message tells the model to ignore conflicts, but a longitudinal study is needed to see whether the default scope is right.

\subsection{What ``recursive'' means here}
The loop is recursive in the sense used by the survey: the state used by later runs is revised using the outcomes of earlier runs, and the update mechanism remains available for further revisions. It is not recursive in the stronger sense of revising the update mechanism itself. The diagnosis heuristics, admission policy, lint rules, and oracle are fixed code. This is a deliberate limit. A system that can revise its own admission policy can lower its own bar; the survey's fixed-evaluator requirement exists for that reason. Meta-level improvement is future work, and it should require an independent audit rather than an autonomous loop.

\subsection{Practical implications}
For a researcher, the useful behaviors are modest but repeated: resolve a DOI before citing, run \texttt{ah\_sweep} after a fact change, audit statistics before reporting, check a funder's policy before drafting aims, and load the right skill before a large analysis. These are not breakthroughs; they are the hygiene that a long-lived assistant should learn. The value of the system is that it learns them from the user's own work, in the user's own vocabulary, without sending the manuscript anywhere and without changing the checks that keep the work honest.

For a harness designer, the lesson is that the editable surface is a design variable. A plugin can expose a rich tool surface and a fixed oracle, and still leave a small, useful space for adaptation. The space does not have to be code. A local lesson store with a visible lifecycle can capture a meaningful part of the adaptation benefit while keeping the evaluation boundary intact.
